\chapter{TOP-D 目标教师的退化与性质}
\label{app:topd}
\section{参数对应与精确可达性}
本附录比较 TOP-D 的目标教师及相关分布性质，依据其 v1 第 3.1 节、定理 4.2、4.4 和附录 C.1、C.2\footfullcite{topd2026}。固定同一前缀，TOP-D 的教师混合权重记为 $\alpha$，教师为 $T$，构造混合时的学生记为 $q$。
\begin{theorem}[TOP-D 目标教师的精确退化]
固定同一条件分布空间的 $q,T>0$。取 $\kappa=1,\eta=\alpha\in(0,1)$ 时，镜像目标教师精确等于 TOP-D 的算术教师 $(1-\alpha)q+\alpha T$，且 $Z=1$。该分布唯一最小化加权正向 KL，超额损失为 $\KL(m\Vert p)$。
\end{theorem}
\begin{proof}
\begin{equation}
\boxed{\kappa=1,\qquad \eta=\alpha,\qquad
m_{1,\eta}=(1-\eta)q+\eta T.}
\label{eq:topd-match}
\end{equation}
这在目标分布上是精确等式，无需小步近似，也无需对两分布接近作假设。两者必须作用于同一个条件分布空间。

目标教师由正文的算术端点最优化问题产生：
\begin{equation}
m_{1,\eta}=\argmin_{p\in\Delta}
\{(1-\eta)\KL(q\Vert p)+\eta\KL(T\Vert p)\}.
\end{equation}
其概率梯度为 $-[(1-\eta)q_a+\eta T_a]/p_a$。约束乘子使 $p_a$ 与分子成比例；分子总和已经为一，所以 $p=m$。完整误差恒等式是
\begin{equation}
F(p)-F(m)=\KL(m\Vert p).
\end{equation}
因此镜像解与该重心解都等于 TOP-D 的目标教师。
\end{proof}

\begin{corollary}[算术教师的总体可达性]
在完整支持、精确期望和目标可实现时，MDPO、Part、BPO 的总体最优分布都等于算术目标教师。Mean 对非恒定信号一般具有不同的最优分布。
\end{corollary}
\begin{proof}
将上一退化定理的算术 $m$ 代入命题~\ref{prop:inner-roots}，MDPO 的 KL gap、Part 的零残差和 BPO 的常数残差条件分别唯一确定 $p=m$。同一命题的 Mean KKT 论证排除非恒定信号下的 Gibbs 解。可达性指总体最优分布；有限参数更新是否到达，应由拟合误差证书判断。
\end{proof}

\section{监督信号、中心与教师目标}
\begin{proposition}[算术端点的信号、中心与教师目标]
取 $\kappa=1$，则
\begin{equation}
f=\log(1-\eta+\eta T/q),\quad A=f/\eta,\quad K_T=\KL(q\Vert m)/\eta.
\end{equation}
本文反推的教师目标为 $J_{1,\eta}(p;T)=\JS_\eta(p,T)/[\eta(1-\eta)]$，其在 $q$ 处步长为 $\eta$ 的负熵镜像解等于算术教师。
\end{proposition}
\begin{proof}
算术端点的配分函数恒为一，故
\begin{equation}
f_{1,\eta}=\log(m/q)=\log(1-\eta+\eta T/q),\qquad
A_{1,\eta}=f_{1,\eta}/\eta,\qquad
K_T=\KL(q\Vert m)/\eta.
\label{eq:topd-signal}
\end{equation}
第一项恢复 TOP-D 式（5）的未缩放信号。第二项是本文内层归一化使用的信号，比较方差时需保留 $1/\eta^2$。

本文由逆向积分定义的教师目标在此端点为
\begin{equation}
J_{1,\eta}(p;T)=
\frac{(1-\eta)\KL(p\Vert M_p)+\eta\KL(T\Vert M_p)}
{\eta(1-\eta)},\qquad M_p=(1-\eta)p+\eta T.
\end{equation}
推导已在附录~\ref{app:objective} 给出。这里 $M_p$ 随候选学生 $p$ 变化，用于定义同一个教师目标；本轮目标 $m=M_q$ 在内层冻结。对 $J$ 在 $q$ 处执行步长 $\eta$ 的负熵镜像下降，恰好得到算术目标。这是本文对该目标教师的变分解释，不把它与其他同名目标混用。
\end{proof}

\section{覆盖与散度收缩}
\begin{proposition}[算术教师的覆盖与信号下界]
设 $m=(1-\eta)q+\eta T$、$0<\eta<1$，则 $m_a\ge(1-\eta)q_a$、$m_a\ge\eta T_a$，并有
\begin{equation}
f\ge\log(1-\eta),\quad
\KL(q\Vert m)\le-\log(1-\eta),\quad \KL(T\Vert m)\le-\log\eta.
\end{equation}
\end{proposition}
\begin{proof}
由算术混合逐坐标有 $m_a\ge(1-\eta)q_a$、$m_a\ge\eta T_a$，因此
\begin{equation}
f_{1,\eta}\ge\log(1-\eta),\quad
\KL(q\Vert m)\le-\log(1-\eta),\quad
\KL(T\Vert m)\le-\log\eta.
\end{equation}
前者控制教师极小概率造成的负监督尾部，后两者控制到两端分布的覆盖代价。
\end{proof}

\begin{proposition}[算术教师的精确收缩与闭式外层]
固定正教师和 $0<\eta<1$，精确算术更新满足 $\|m-T\|_1=(1-\eta)\|q-T\|_1$。反复更新的闭式为 $q_k=T+(1-\eta)^k(q_0-T)$。
\end{proposition}
\begin{proof}
对于精确更新，
\begin{equation}
m-T=(1-\eta)(q-T),\qquad
\|m-T\|_1=(1-\eta)\|q-T\|_1.
\end{equation}
反复精确更新的闭式为 $q_k=T+(1-\eta)^k(q_0-T)$。所以原文定理 4.4 所用的算术收缩结构在这个条件分布空间直接恢复。
\end{proof}

\begin{proposition}[算术端点的凸散度收缩]
对任意凸函数 $\varphi$、$\varphi(1)=0$，
\begin{equation}
D_\varphi(m\Vert T)\le(1-\eta)D_\varphi(q\Vert T).
\label{eq:topd-all-f}
\end{equation}
\end{proposition}
\begin{proof}
逐坐标有 $m_a/T_a=(1-\eta)q_a/T_a+\eta$。凸性给
\begin{equation}
\varphi(m_a/T_a)
\le(1-\eta)\varphi(q_a/T_a)+\eta\varphi(1).
\end{equation}
乘以 $T_a$ 并求和，第二项为零。取 $\varphi(x)=x\log x$、$-\log x$、$(\sqrt x-1)^2$、$|x-1|/2$，分别得到两个 KL 方向、平方 Hellinger 和 TV 的界。
\end{proof}

式~\eqref{eq:topd-all-f} 是算术映射自身的凸性结论。一般幂均值保留所有凸生成散度的非增性，但不自动具有同一个 $(1-\eta)$ 全局因子；平方根点的 Hellinger 因子由它的球面结构另行给出。

\section{有限拟合的误差累积}
\begin{theorem}[算术外层的拟合误差累积]
若 $\epsilon_k=\|q_{k+1}-m(q_k,T)\|_1$，则
\begin{equation}
\|q_K-T\|_1\le(1-\eta)^K\|q_0-T\|_1+\sum_{k<K}(1-\eta)^{K-1-k}\epsilon_k.
\end{equation}
若 $\limsup\epsilon_k\le\epsilon_\infty$，则 $\limsup\|q_K-T\|_1\le\epsilon_\infty/\eta$。固定测度的条件距离平均满足同一界。
\end{theorem}
\begin{proof}
令 $\epsilon_k=\|q_{k+1}-m(q_k,T)\|_1$。三角不等式给
\begin{equation}
\|q_{k+1}-T\|_1\le(1-\eta)\|q_k-T\|_1+\epsilon_k.
\end{equation}
归纳展开：
\begin{equation}
\|q_K-T\|_1\le(1-\eta)^K\|q_0-T\|_1
+\sum_{k=0}^{K-1}(1-\eta)^{K-1-k}\epsilon_k.
\label{eq:topd-errors}
\end{equation}
若 $\limsup\epsilon_k\le\epsilon_\infty$，对任意 $\delta>0$ 取足够晚的 $k_0$ 使 $\epsilon_k\le\epsilon_\infty+\delta$。有限前缀和趋零，后段几何和不超过 $(\epsilon_\infty+\delta)/\eta$。令 $\delta\downarrow0$，得到 $\limsup\|q_K-T\|_1\le\epsilon_\infty/\eta$。这恢复原文误差累积与稳态上界的形式。

固定前缀测度下可以逐条件取平均，再用同一个线性递推。如果用 MDPO 的超额损失 $\delta_k$，Pinsker 给 $\epsilon_k\le\sqrt{2\eta\delta_k}$；若用 BPO，则先应用附录~\ref{app:dynamics} 的跨度证书。本文在这里补上的是不同 inner loss 到外层输出误差的接口。
\end{proof}

\section{方差定理的尺度对应}
\begin{theorem}[算术监督的方差界与小步尺度]
固定 $0<\eta<1$，$f=\log(1-\eta+\eta T/q)$，且 $\|\psi\|\le S$。令 $C_\star=\sup_{u>0}\log^2(1+u)/u$，则
\begin{equation}
\Var_q(f\psi)\le S^2[\log^2(1-\eta)+C_\star\eta].
\end{equation}
对正分布类取上确界，$\sup\Var_qf=\Theta(\eta)$，而 $A=f/\eta$ 满足 $\sup\Var_qA=\Theta(\eta^{-1})$ 当 $\eta\downarrow0$。
\end{theorem}
\begin{proof}
TOP-D 定理 4.2 在有限词表和 score 范数不超过 $S$ 的条件下，给未缩放信号的梯度方差上界
\begin{equation}
S^2|\mathcal V|\max\{\log^2(1-\eta),C_\star\eta\}.
\end{equation}
在同一条件下，利用概率归一化可以得到无词表大小因子的界。对 $f=\log(1-\eta+\eta\rho)$，将 $\rho<1$ 与 $\rho\ge1$ 分开：
\begin{align}
\E_q[f^2\mathbf1_{\rho<1}]&\le\log^2(1-\eta),\\
\E_q[f^2\mathbf1_{\rho\ge1}]
&\le C_\star\eta\E_q[\rho\mathbf1_{\rho\ge1}]
\le C_\star\eta.
\end{align}
其中 $C_\star=\sup_{u>0}\log^2(1+u)/u$，和原文单变量上界使用同一个常数。因此
\begin{equation}
\Var_q(f\psi)\le S^2[\log^2(1-\eta)+C_\star\eta].
\end{equation}
这个推导控制条件单动作梯度；前缀加权后只要使用同一 score 上界，二阶矩界仍成立。

未缩放信号的最坏二阶矩在 $\eta\downarrow0$ 时为 $\Theta(\eta)$。上界由前式，下界取 $q_1=\eta/2,T_1=1/2$，稀有动作 $f_1=\log(2-\eta)$，贡献至少 $\eta\log^2(2-\eta)/2$。对二元分布，中心化方差 $q_1q_2(f_1-f_2)^2$ 也有同阶，因为另一项 $f_2=O(\eta)$。

所以本文标准化监督满足
\begin{equation}
\sup_{q,T}\Var_qA_{1,\eta}=\Theta(\eta^{-1}).
\end{equation}
这与未缩放信号随小 $\eta$ 衰减相容：两者相差 $\eta^2$。它也说明算术端点虽有固定强度方差界，却不属于统一小步标准化方差有界的 $0<\kappa\le1/2$ 区间。
\end{proof}

\section{比较结论的作用范围}
\begin{center}
\small
\begin{tabular}{p{38mm}p{93mm}}
\toprule
对象 & 精确对应或进一步性质\\
\midrule
目标教师 & $\kappa=1,\eta=\alpha$ 时完全一致，$Z=1$\\
未缩放信号 & $f=\log(1-\eta+\eta T/q)$ 完全一致\\
最优化来源 & 加权正向 KL 重心，超额损失为 $\KL(m\Vert p)$\\
教师目标解释 & 本文的 $J_{1,\eta}$ 是归一化加权 JS，其熵镜像步等于混合教师\\
精确与近似外层 & 算术 $L_1$ 收缩、几何误差累积、$\epsilon_\infty/\eta$ 稳态界\\
其他教师散度 & 凸性给任意凸生成散度的 $(1-\eta)$ 上界\\
监督方差 & 未缩放 $f$ 与标准化 $A=f/\eta$ 分开比较\\
\bottomrule
\end{tabular}
\end{center}

上述比较针对相同的条件分布空间。若直接在完整序列空间定义算术混合，同一代数也成立；逐 token 固定权重混合一般不等于完整序列混合，不能把两个空间的距离直接认作同一指标。这里保留目标教师与定理层面的对应。
